%% ---------------------------------------------------------------------
\section{Función de transferencia}
\label{sec:po-ft}
%% ---------------------------------------------------------------------

Las herramientas anteriores nos llevan al concepto central de todo el curso.

\begin{definicion}[label=def:po-ft]{Función de transferencia}
La \clave{función de transferencia} de un sistema LTI es el cociente entre la transformada de Laplace de la salida y la de la entrada, \emph{con condiciones iniciales nulas}:
\begin{equation}
  T(s)=\left.\frac{\Lap{\text{salida}}}{\Lap{\text{entrada}}}\right|_{\text{CI}=0}=\frac{Y(s)}{U(s)}.
\end{equation}
\end{definicion}

Si el sistema se describe con la ecuación entrada-salida
\begin{equation}
  \frac{\dif^n y}{\dif t^n}+a_1\frac{\dif^{n-1}y}{\dif t^{n-1}}+\cdots+a_n y
  =b_0\frac{\dif^m u}{\dif t^m}+\cdots+b_{m-1}\frac{\dif u}{\dif t}+b_m u,
  \label{eq:po-inout}
\end{equation}
la linealidad y la diferenciación con condiciones iniciales nulas dan
\begin{equation}
  T(s)=\frac{Y(s)}{U(s)}=\frac{b_0s^m+b_1s^{m-1}+\cdots+b_m}{s^n+a_1s^{n-1}+\cdots+a_n}.
  \label{eq:po-ft}
\end{equation}

Algunos nombres que usaremos todo el tiempo:
\begin{itemize}
  \item $T(s)$ es \clave{propia} si $n\geq m$, y \clave{estrictamente propia} si $n>m$. Los sistemas físicos son propios: una función impropia derivaría la entrada y amplificaría sin límite el ruido de alta frecuencia.
  \item El \clave{grado relativo} es $n-m$. El \clave{orden} del sistema es $n$.
  \item El polinomio del denominador igualado a cero, $s^n+a_1s^{n-1}+\cdots+a_n=0$, es la \clave{ecuación característica}. Sus raíces son los polos de $T(s)$: los modos naturales del sistema.
\end{itemize}

La función de transferencia permite describir el sistema con un bloque (\cref{fig:po-bloqueft}): la señal de salida es el producto $Y(s)=T(s)\,U(s)$. En el tiempo, por la propiedad de convolución, eso es $y(t)=\int_0^t g(\lambda)\,u(t-\lambda)\,\dif\lambda$, donde $g(t)=\iLap{T(s)}$ es la \clave{respuesta al impulso}: el producto en $s$ es mucho más cómodo que la integral en $t$.

\begin{figure}[ht]
\centering
\includegraphics[width=0.62\linewidth]{cap02/bloque_ft}
\caption{Un bloque de flujo de información: la función de transferencia es un mapeo lineal entre la transformada de la entrada y la de la salida. Las reglas para combinar bloques están en el \cref{cap:so}.}
\label{fig:po-bloqueft}
\end{figure}

\begin{cuidado}[La función de transferencia no ve las condiciones iniciales]
$T(s)$ describe la respuesta \emph{forzada} desde el reposo. Si el sistema no arranca en reposo, la respuesta total es la forzada más la libre (la que producen las condiciones iniciales, como en el ejemplo~\ref{ej:po-motor}). Por eso trabajamos en variables desviación alrededor de un equilibrio: allí las condiciones iniciales son cero.
\end{cuidado}

\begin{pausa}
Para $\dddot y+2\ddot y+3\dot y+4y=\dddot u+5\dot u+6u$, escriba $T(s)$. ¿Es propia? ¿Cuál es su grado relativo?
\respuesta{$T(s)=\dfrac{s^3+5s+6}{s^3+2s^2+3s+4}$. Es propia pero no estrictamente propia ($n=m=3$): grado relativo $0$. Una parte de la entrada pasa directamente a la salida.}
\end{pausa}
